% Incorporación TER-01--04, TER-06--07. TER-05 conserva su desarrollo en
% 09_area, 09b_hojas_memoria, 09c_entropia_sectorial y la composición posterior.
% Fuentes: BOLTZMANN_DESARROLLO_OPERATORIO_20260927/{retorno,normalizacion,
% incidencia,seleccion,accion_horizonte_20260927}; FOCK_PARIDAD_RAMAS_Y_SELECTOR;
% REVISION_LIFT_POSITIVO_Y_ESPECTRO; ADDENDUM_TRANSPORTE_TERMICO_Y_SECCIONES.
% Edición aditiva: no sustituye el núcleo ni las demostraciones anteriores.
% Recibo causal heredado: APP conserva suma, producto, residuos y cocientes.
% TRIT conserva régimen y orientación en el calendario recibido.
% TPK transporta ruta, carry, frontera y memoria en la conexión nonádica.
% El estado enriquecido prolonga el registro al retornar la fase.
% La estructura discreta conjunta del continuo es el dominio heredado.
% Salida HMT examinada: capacidad, portador graduado y lector marcado.
% La realización espectral convencional es posterior y conserva las reglas declaradas.
\section{Capacité spectrale, mémoire énergétique et transport thermique}
\label{sec:ii-ter27-antecedentes}

La construction suivante réunit la capacité des prolongements, le
renouvellement des deux feuillets et leurs réalisations spectrales. Elle vise à
déterminer ce que conserve une lecture thermique lorsque la phase revient et que le
registre continue de croître. Les opérations proviennent de l’état enrichi
APP--TRIT--TPK : le catalogue des signatures, le calendrier, l’orientation et la
mémoire précèdent les traces qui seront évaluées. Les espaces linéaires,
les compressions et les préparations ajoutés à ces antécédents
sont expressément définis. Chaque résultat conserve ainsi ses conditions de
construction et sa fonction dans l’argument.

\subsection{Capacité, vacance et renouvellement conjoint}
\label{sec:ii-ter27-capacidad}

La comparaison entre les $729$ mots de six trits et les $1000$
positions d’un bloc décimal définit, avec l’origine déjà fixée dans
\ref{vac-capacidad-trit},
\begin{equation}
 C_t=\lfloor(t+1)\rho\rfloor,
 \qquad \rho=\log_{1000}729=\log_{10}9,
 \qquad z_t=1-(C_t-C_{t-1}).
 \label{eq:ii-ter27-capacidad}
\end{equation}
L’inégalité $0<\rho<1$ implique $z_t\in\{0,1\}$. La capacité compte
des blocs dans cette comparaison ; son évaluation exacte s’effectue au moyen de
$10^{C_t}\leq9^{t+1}<10^{C_t+1}$. Pour $1\leq t\leq20$, on obtient
$C_t=t$, tandis que $C_{21}=20$. En effet, pour $t\leq20$, l’inégalité
inférieure équivaut à $9(9/10)^t\geq1$, et il suffit de la vérifier en $t=20$ ;
en $t=21$, cette inégalité s’inverse. Les deux comparaisons restantes
situent $9^{22}$ entre $10^{20}$ et $10^{21}$. La première vacance positive
est donc $p_1=21$, avec
\[
 C_{20}=C_{21}=20,\qquad C_{22}=21,\quad C_{23}=22,\quad C_{24}=23.
\]
Les vingt accroissements de capacité forment vingt étiquettes effectives.
L’indice de coupure $24$, les $25$ positions de $0$ à $24$ et la capacité
$23$ sont des registres de types distincts.

Le catalogue conjoint des émissions est
$\Omega=\mathcal S_+\times\mathcal S_\times$, de cardinaux $18$ et
$26$. Sur $\ell^2(\Omega)$ avec comptage uniforme, les renouvellements de
la coordonnée active sont
\[
 E_+=\frac{\mathbf J_{18}}{18}\otimes I_{26},\qquad
 E_\times=I_{18}\otimes\frac{\mathbf J_{26}}{26}.
\]
Ce sont des projecteurs orthogonaux qui commutent. Leur produit est le projecteur
constant $\Pi$, dont la matrice a toutes ses entrées égales à $1/468$.
La lecture réduite de phase utilise le régime d’origine :
\[
 K_{\rm ph}((u,r),(v,r+1))=P_{\tau(r)}(u,v),
 \qquad (P_+,P_-,P_0)=(E_+,E_\times,I),
 \quad \tau(r)=M_{\rm ph}(1+r).
\]
Cette opération est postérieure à la dynamique enrichie et à son holonomie,
dans l’ordre $U_t\longrightarrow\Gamma_9\longrightarrow K_{\rm ph}$.

\begin{proposition}[Renouvellement des deux feuillets sur l’horloge de capacité]
\label{prop:ii-ter27-renovacion}
Évaluons le lecteur précédent au moyen de $\mathcal R_t=K_{\rm ph}^{a_t}$,
où $a_t=C_{t+1}-C_t$. Depuis la première vacance, le premier instant où
sa sortie matricielle est indépendante de l’émission initiale est $24$.
Les rangs successifs sont $468,26,1$ ; dans l’orientation conjuguée, ce sont
$468,18,1$.
\end{proposition}
\begin{proof}
Les capacités d’origine $20,21,22$ correspondent aux régimes
$0,+1,-1$. Les avancées $21\to22\to23\to24$ appliquent donc
$I,E_+,E_\times$, et leurs produits cumulés sont $I,E_+,\Pi$.
Le rang d’un produit tensoriel est le produit de leurs rangs :
$\operatorname{rank}E_+=26$, $\operatorname{rank}E_\times=18$ et
$\operatorname{rank}\Pi=1$. Les lignes de $I$ et $E_+$ dépendent de
l’émission initiale ; celles de $\Pi$ sont identiques. Cela prouve la minimalité.
La séquence conjuguée applique $I,E_\times,E_+$ et a le même produit
final. La composition par capacité conserve le temps chronologique :
\[
 (t,C_t,v_t)\mapsto(t+1,C_t+a_t,v_t+1-a_t),\qquad v_t=t-C_t.
\]
L’identité $t=C_t+v_t$ reste valable également lorsque $a_t=0$.
\end{proof}

La descente de la phase modulo neuf au régime modulo trois est exacte
pour ce lecteur. Pour $Q(u,r)=(u,[r]_3)$,
\[
 \sum_{s:[s]_3=b}K_{\rm ph}((u,r),(v,s))
 =\mathbf1_{b=[r+1]_3}P_{\tau(r)}(u,v).
\]
La somme dépend uniquement de $Q(u,r)$, ce qui prouve l’agrégation forte.
En revanche, le relèvement conserve $C=9m+3b+a$. Le passage $20\to23$
transforme $(m,b,a)=(2,0,2)$ en $(2,1,2)$ : $a$ revient et $b$ avance.
La phase positive modulo neuf passe de $3$ à $6$. Le renouvellement uniforme
est une sortie matricielle de rang un ; le chemin et le microétat antérieur sont
conservés dans l’archive enrichie, non dans cette seule sortie matricielle.
L’information d’un régime uniforme est $\log3$ ; celle des émissions
renouvelées est $\log468$ ; le taux d’une rotation déterministe est nul.
Ces trois lectures correspondent à des ensembles et à des opérations différents.

\subsection{Retours par capacité et valuation des cylindres}
\label{sec:ii-ter27-cociclo}

Définissons le retour strict du régime par
\[
 R_3(s)=\min\{t>s:C_t>C_s,\ C_t\equiv C_s\pmod3\}.
\]
La suite $C_t$ est non bornée et ses accroissements valent zéro ou un.
Elle visite donc $C_s+3$ avant toute autre valeur positive congruente
à $C_s$ ; ainsi, $C_{R_3(s)}=C_s+3$. À partir de $s_0=24$, on obtient la
localisation corrigée
\begin{equation}
 s_j=\min\{t\geq24:C_t=23+3j\},\qquad C_{s_j}=23+3j.
 \label{eq:ii-ter27-localizacion}
\end{equation}
Les premiers temps sont $24,27,30$. Les intervalles suivants peuvent
contenir des vacances ; l’augmentation de capacité reste exactement de trois.

Pour expliquer la valuation, décomposons chaque bloc comme
$d_r=100a_{r,1}+10a_{r,2}+a_{r,3}$, avec $0\leq a_{r,k}\leq9$.
Un préfixe de $c$ blocs se transporte vers les trois entiers
$X_k=\sum_{r=1}^c a_{r,k}10^{c-r}$. La division euclidienne fournit
une bijection ; son inverse regroupe les chiffres de chaque position. La
troncature d’un bloc devient la troncature d’un chiffre
dans chaque coordonnée. La somme avec retenue se transporte au moyen de cette
bijection complète, en conservant ses quotients.

Le cylindre scalaire de longueur $1000^{-c}$ correspond au cube semi-ouvert
\[
 \prod_{k=1}^3[X_k10^{-c},(X_k+1)10^{-c}),\qquad
 J_c=10^{-c}I_3,\qquad \det J_c=1000^{-c}.
\]
L’égalité exprime une correspondance entre les mesures des cylindres. Entre
deux capacités, $J_{c',c}=10^{-(c'-c)}I_3$ satisfait
$J_{c'',c'}J_{c',c}=J_{c'',c}$. La similitude positive isotrope est
déterminée par son déterminant, car $s^3=1000^{-(c'-c)}$ a une
unique racine positive. Aux retours marqués,
\begin{equation}
 J_{C_{s_j}}=10^{-23}1000^{-j}I_3.
 \label{eq:ii-ter27-cociclo}
\end{equation}
Le facteur $10^{-23}$ est l’échelle d’une coordonnée du cube ; la mesure
du cylindre scalaire à cette capacité est $10^{-69}$. La distinction conserve
le degré géométrique du lecteur. Un extracteur fixé en $s_0$, suivi d’une
troncature compatible, satisfait $E_n\pi_{m,n}=E_m$ pour $m\geq n\geq
s_0$ ; le prolongement conserve ainsi le registre marqué sans recalculer ses
rangs à partir d’une coupure ultérieure.

\subsection{Porteurs de relations et réalisation par échange}
\label{sec:ii-ter27-portador}

Soient $U_m$ l’espace des étiquettes de capacité et $P_p$ celui des positions
du préfixe. L’espérance diagonale sur $\operatorname{End}(U_m)$ est
$E_{\rm diag}(X)=\sum_iE_{ii}XE_{ii}$. Son complément orthogonal a pour
base $E_{ij}$, $i\ne j$, et pour dimension $m(m-1)$. Le second espace de
registres est $\operatorname{End}(P_p)\oplus A_{p-1}$, où $A_{p-1}$
conserve les transitions adjacentes comme observations typées. Sa
dimension est $p^2+p-1$. Avec $m=20$, $p=25$, les rangs sont $380$ et $649$.
Ce choix d’espaces et de coupure définit une réalisation du lecteur ;
les identités suivantes démontrent ses opérations sans attribuer au
comptage, à lui seul, la sélection physique de ce domaine.

Le renversement basé agit par
$E_{ij}\mapsto E_{m-1-i,m-1-j}$ et, dans le second espace, par
$E_{ij}\mapsto E_{p-1-i,p-1-j}$, $a_j\mapsto a_{p-2-j}$.
Pour $m=20$, il y a $190$ orbites à deux éléments. Pour $p=25$, il existe un
unique élément matriciel fixe, $E_{12,12}$ ; il reste $312$ orbites
matricielles et $12$ orbites de transitions. On obtient une droite fixe et
$324$ représentations régulières de $C_2$, avec des espaces propres de dimensions
$325$ et $324$. C’est l’antécédent de $649=1+2\cdot324$ ; il conserve
l’involution des registres, il ne détermine pas de graduation énergétique.

Une seconde réalisation utilise le même échange $J$ de deux
orientations et conserve les termes tensoriels croisés. Soient $W,V$
des espaces de dimensions $m,n$. Dans le secteur extérieur de $W_+\oplus W_-$
et dans les oscillateurs symétriques de $V_+\oplus V_-$, les caractères fixes sont
\begin{align}
 \chi_F(q)&=\tfrac12\bigl((1+q)^{2m}+(1-q^2)^m\bigr),\\
 \chi_B(q)&=\tfrac12\bigl(P_n(q)^2+P_n(q^2)\bigr),
 &P_n(q)&=\prod_{r\geq1}(1-q^r)^{-n}.
 \label{eq:ii-ter27-caracteres}
\end{align}
\begin{proof}
À une particule, les espaces propres $J=\pm1$ ont tous deux pour dimension $m$.
La trace avec insertion dans l’algèbre extérieure est
$(1+q)^m(1-q)^m=(1-q^2)^m$. Le projecteur $(I+J)/2$ fait la moyenne de cette trace
et du caractère ordinaire. Le coefficient de degré deux est
$\bigl(\binom{2m}{2}-m\bigr)/2=m(m-1)$.
Pour chaque fréquence bosonique, les multiplicités des deux signes sont
$n$, donc la trace avec insertion est
$\prod_r[(1-q^r)(1+q^r)]^{-n}=P_n(q^2)$.
Les coefficients ordinaires aux degrés zéro, un et deux sont
$1,2n,2n^2+3n$, et ceux avec insertion sont $1,0,n$. Leurs moyennes sont
$1,n,n(n+2)$. Chaque coefficient formel dépend d’un nombre fini de
facteurs, ce qui justifie le produit sans hypothèse analytique supplémentaire.
\end{proof}
Le terme supplémentaire $n$ à l’énergie deux provient des créations de
fréquence deux. En effet,
$M(V)_2=V(-2)\oplus\operatorname{Sym}^2(V(-1))$ ; pour le rang $24$,
sa dimension est $24+300=324$. Dans les deux orientations échangées,
le sous-espace fixe bosonique cumulé jusqu’à l’énergie deux a pour dimension
$1+n+n(n+2)=(n+1)^2+n$, ce qui donne $649$ pour $n=24$.
L’identification linéaire avec l’espace précédent de matrices et de transitions
nécessite de transporter leurs bases ; l’égalité des dimensions n’identifie pas
automatiquement leurs hamiltoniens. La coupure extérieure homogène et la coupure
bosonique cumulée restent déclarées comme parties de ce lecteur.

\subsection{Germes multiplicatifs et compensation exacte des fibres}
\label{sec:ii-ter27-fibras}

Le support des germes multiplicatifs est
$X_\times=(\mathbb Z/9\mathbb Z)^2\times\{N,E,S,O\}$, de cardinal
$324$. L’émetteur à six fenêtres $e:X_\times\to\mathcal S_\times$
du noyau produit $26$ signatures. Ses fibres ont pour cardinaux $8$ pour
$24$ signatures, $24$ pour $555555$ et $108$ pour $000000$ ; la somme est $324$.
L’ensemble temporel $I=\{1,\ldots,24\}$ conserve un type distinct de la
fibre de $555555$, bien que tous deux aient pour cardinal $24$.

On construit l’espace microscopique
\[
 \widetilde T=\{\Omega\}\sqcup\{a_i:i\in I\}
 \sqcup\{b_{i,y}:i\in I,y\in X_\times\},\qquad |\widetilde T|=7801.
\]
La lecture $p(b_{i,y})=b_{i,e(y)}$, qui fixe les autres registres, a pour
image $T$ de cardinal $1+24+24\cdot26=649$. Soit $m_e=|e^{-1}(e)|$.
Assigner une masse $1$ à $\Omega,a_i$ et $1/m_{e(y)}$ à $b_{i,y}$ définit
une mesure $\nu$ dont les fibres par $p$ ont toutes une masse égale à un. Pour
$r>0$, les taux sont $r$ dans les deux directions entre $\Omega$ et $a_i$,
et $r/m_{e(y)}$ de $a_i$ vers $b_{i,y}$, avec un taux inverse $r$.

\begin{proposition}[Transport compensé de l’arbre des signatures]
\label{prop:ii-ter27-arbol}
Le générateur microscopique est réversible pour $\nu$. L’application $p$
induit l’arbre symétrique sur $T$ et le relèvement $J_pg=g\circ p$
est une isométrie qui entrelace les deux générateurs.
\end{proposition}
\begin{proof}
Sur chaque arête terminale, on a le bilan
$\nu(a_i)r/m_{e(y)}=\nu(b_{i,y})r$. De plus,
$\sum_{y:e(y)=e}r/m_e=r$, et le taux inverse a la même valeur depuis
tout représentant. Cela prouve l’agrégation forte. L’identité
$\sum_{x:p(x)=z}\nu(x)=1$ prouve $J_p^*J_p=I$ ; les sommes de taux
prouvent $\widetilde L J_p=J_pL_T$. L’arbre est connexe et a
$648$ arêtes, de sorte que sa mesure stationnaire normalisée est uniforme.
\end{proof}
Le demi-tour du curseur induit une involution des signatures qui conserve
$m_e$. Combinée à $i\mapsto25-i$, elle conserve $\nu$ et les taux ;
l’entrelaceur transporte également cette inversion. Les germes restent
enregistrés dans $\widetilde T$ bien que $T$ ne publie que leur signature.
La sélection d’une marque et le rafraîchissement compensé sont les opérations
de ce lecteur. Son domaine et ses taux restent explicites : le chemin
cardinal d’un curseur et la mise à jour enrichie complète conservent
des opérations supplémentaires. L’arbre relie les niveaux internes du
secteur bosonique ; les communications entre les trois degrés seront construites
dans \ref{sec:ii-ter27-corrientes}.

\subsection{Double graduation et récupération de ses marginales}
\label{sec:ii-ter27-doble}

Le porteur sélectionné pour la comparaison est
\[
 \mathcal D=\mathbb C\Omega_{\rm ext}\oplus F_2^J\oplus B_{\leq2}^J.
\]
L’opérateur $L$ conserve le degré de Fock ; $N$ assigne les degrés externes
$0,1,2$ aux trois facteurs directs. Avec $m=20,n=24$, sa décomposition conjointe est
\begin{equation}
 (N,L;d)=(0,0;1),(1,2;380),(2,0;1),(2,1;24),(2,2;624).
 \label{eq:ii-ter27-celdas}
\end{equation}
Les deux vides sont orthogonaux et conservent leur provenance. La fonction
bivariée et ses spécialisations sont
\begin{align}
 \mathcal Z(u,v)&=1+380uv^2+u^2(1+24v+624v^2),\\
 Z_N(q)&=\mathcal Z(q,1)=1+380q+649q^2,\\
 Z_L(q)&=\mathcal Z(1,q)=2+24q+1004q^2.
 \label{eq:ii-ter27-particiones}
\end{align}
Chaque monôme résulte de l’application de $u^Nv^L$ à une cellule de
\eqref{eq:ii-ter27-celdas} suivie de la prise de sa trace. En $q=1/1000$, on obtient
respectivement $1.380649$ et $2.025004$. La seconde évaluation réunit deux
vides, $24$ états d’énergie un et $1004$ d’énergie deux.

Les marginales complètes $n_i,l_j$ d’une distribution supportée dans
ces cinq cellules permettent de retrouver leurs masses :
\[
 p_{00}=n_0,\quad p_{12}=n_1,\quad p_{21}=l_1,
 \quad p_{20}=l_0-n_0,\quad p_{22}=l_2-n_1.
\]
La somme par lignes et colonnes prouve la restitution dès lors que les
différences sont non négatives et que les totaux coïncident. Le support est un
arbre biparti ; la restitution concerne ses cinq masses, tandis que
les cohérences et les indices internes restent dans l’état complet.
Deux valeurs de partition isolées ne remplacent pas ces marginales. Comme
contrôle, $Z_N-Z_L=(q-1)(1-355q)$ : en $q=1/355$, les deux traces coïncident,
mais les probabilités du vide interne sont $1/261574$ et
$126025/261574$. L’égalité du scalaire conserve ici des états différents.

L’inversion basée du lecteur de préfixes permet également d’évaluer
l’information d’orientation omise par son quotient. Il y a $190$ paires au
degré un et $324$ au degré deux ; chaque paire équipondérée apporte $\log2$
fois sa masse. Par conséquent,
\[
 I_{\rm or}(q)=\frac{380q+648q^2}{Z_N(q)}\log2.
\]
La restitution de l’étiquette d’orientation retrouve cette information.
Le quotient et sa lecture entropique sont distingués d’une opération
physique d’effacement.

\subsection{Dilatation énergétique positive de la mémoire}
\label{sec:ii-ter27-memoria}

Comme $N,L$ commutent, $Q_c=N-L$ prend les valeurs $0,-1,2,1,0$ dans les
cellules précédentes. Soit $J_m=\operatorname{diag}(0,1,2,3)$ et $g>0$.
Définissons l’isométrie
\begin{equation}
 V|n,l,i\rangle=|n,l,i\rangle\otimes|n-l+1\rangle,
 \quad H_F=gL,\quad H_m=gJ_m.
 \label{eq:ii-ter27-lift}
\end{equation}
Les niveaux auxiliaires sont $1,0,3,2,1$. La première coordonnée conserve
la base complète, donc $V^*V=I$. Sur chaque vecteur,
\begin{equation}
 (H_F\otimes I+I\otimes H_m)V=Vg(N+I),
 \label{eq:ii-ter27-energia-memoria}
\end{equation}
parce que $l+(n-l+1)=n+1$. L’énergie auxiliaire est non négative et les
énergies totales sont $g,2g,3g,3g,3g$. L’image est réductrice pour ce
hamiltonien diagonal. Pour $q=e^{-\beta g}$, le calcul fonctionnel donne
\[
 V^*e^{-\beta H_{\rm tot}}V=q^{N+I},\qquad
 \frac{V^*e^{-\beta H_{\rm tot}}V}
 {\operatorname{Tr}(V^*e^{-\beta H_{\rm tot}}V)}=\frac{q^N}{Z_N(q)}.
\]
Le filtre $M_q=q^{(Q_c+I)/2}$ satisfait
$V^*(I\otimes q^{J_m})V=M_q^*M_q$ et
$q^L M_q^*M_q=q^{N+I}$. Pour $0<q<1$, il est contractant et réalise le
passage conditionné de $\rho_L=q^L/Z_L$ à $\rho_N=q^N/Z_N$ avec la
probabilité $qZ_N/Z_L$. Enregistrer son résultat conserve également la
branche complémentaire. L’isométrie $V$ conserve les cohérences ; la trace
partielle sur la mémoire constitue une autre opération qui élimine uniquement les
cohérences entre différentes valeurs de $Q_c$.
Le décalage $+1$ est le plus petit entier qui rend non négatifs tous
les niveaux auxiliaires. Un décalage commun plus grand conserve l’état
normalisé et change l’origine énergétique enregistrée.

\subsection{Boucle de mémoire et spectre décimal}
\label{sec:ii-ter27-espectro}

La réalisation gaussienne reçoit la décomposition collective $8:1$ des
neuf copies et la boucle unitaire de mémoire. Les deux modes sont préparés
dans la référence gaussienne pure centrée de covariance conjointe $I\oplus I$,
dans les mêmes coordonnées symplectiques que celles où sont écrites les matrices
suivantes :
\[
 Q_8=\frac13\begin{pmatrix}\sqrt8 I&-I\\ I&\sqrt8 I\end{pmatrix},
 \quad H_1=\begin{pmatrix}2&1\\1&1\end{pmatrix},
 \quad U_{H_1}=Q_8^T\operatorname{diag}(I,H_1)Q_8.
\]
Ces matrices reçoivent leur préparation et leur boucle du transport du
système et de la mémoire ; le calcul thermique s’effectue ensuite. Le bloc
visible de $U_{H_1}U_{H_1}^T$ est
\[
 W=\frac{8I+H_1H_1^T}{9}
 =\frac19\begin{pmatrix}13&3\\3&10\end{pmatrix},
 \qquad \det W=\frac{121}{81},\quad \nu=\frac{11}{9}.
\]
En dimension deux, $S=\sqrt\nu W^{-1/2}$ a pour déterminant un et est
symplectique, avec $SWS^T=\nu I$. Sa réalisation métaplectique réduit
l’état gaussien à cette covariance isotrope. L’état géométrique
$(1-r)\sum_{j\geq0}r^j|j\rangle\langle j|$ a pour occupation moyenne
$r/(1-r)$ et pour covariance $(1+r)/(1-r)I$. L’unicité de l’état gaussien
centré par sa covariance implique
\begin{equation}
 r=\frac{\nu-1}{\nu+1}=\frac1{10},\qquad
 \rho_r=(1-r)\sum_{j\geq0}r^j|j\rangle\langle j|.
 \label{eq:ii-ter27-williamson}
\end{equation}
La raison décimale résulte de la boucle et de la préparation indiquées. Par
exemple, $H_0=I$ produirait $\nu=1$ et $r=0$ ; le résultat conserve
ainsi sa dépendance effective à l’égard de l’opérateur de mémoire.

La division euclidienne $j=3w+b$, $0\leq b<3$, définit un opérateur unitaire
$U_3|j\rangle=|w,b\rangle$ d’inverse $|w,b\rangle\mapsto|3w+b\rangle$.
La factorisation $1-r^3=(1-r)(1+r+r^2)$ démontre terme à terme
\begin{equation}
 U_3\rho_rU_3^*=\rho_q\otimes\sigma_r,
 \quad q=r^3=\frac1{1000},\quad
 \sigma_r=\frac{\operatorname{diag}(100,10,1)}{111}.
 \label{eq:ii-ter27-agrupamiento}
\end{equation}
Le décalage unilatéral se transporte en
\[
 T=I\otimes(|1\rangle\langle0|+|2\rangle\langle1|)
       +S_w\otimes|0\rangle\langle2|,\qquad T^3=S_w\otimes I_3.
\]
L’application à chaque vecteur de base prouve cette identité et montre la
retenue du résidu vers le quotient. On a $T^*T=I$ et
$TT^*=I-|0,0\rangle\langle0,0|$ ; le retour du résidu conserve
l’avancée d’occupation. L’application unitaire transporte également toute purification
en agissant comme $U_3\otimes I$ sur le système étendu.

Pour la section énergétique $H_W=\epsilon_a(j+1/2)$, le transport donne
\[
 U_3H_WU_3^*=\epsilon_a(3N_w+B+\tfrac12 I),
 \qquad B=\operatorname{diag}(0,1,2).
\]
L’identité diagonale conserve le domaine autoadjoint déterminé par
la sommabilité quadratique de l’énergie. À $\beta_W\epsilon_a=\log10$,
le quotient a pour espacement $3\epsilon_a$ et pour raison $q$, en conservant
le même thermomètre. Leurs moyennes satisfont
$\langle N_w\rangle=1/999$, $\langle B\rangle=4/37$ et
$3\langle N_w\rangle+\langle B\rangle=1/9$.
De plus, $s(\rho_r)=s(\rho_q)+s(\sigma_r)$, par la factorisation exacte.

Cette précision fixe $g=3\epsilon_a$ dans la composition de
\eqref{eq:ii-ter27-lift} avec le quotient de Williamson. Si l’on prépare
$\rho_L\otimes\rho_q$ et que l’on enregistre la projection sur l’image de $V$,
\[
 V^*(\rho_L\otimes\rho_q)V=\frac{1-q}{Z_L}q^{N+I}.
\]
La probabilité est $(1-q)qZ_N/Z_L$ et l’état conditionné est $\rho_N$.
Le facteur $\sigma_r$ est conservé au moyen de $V\otimes I_3$ ; l’énergie
totale transportée est
$3\epsilon_a(N+1)+\epsilon_a(B+1/2)$.

\subsection{Section spectrale marquée et intensité complète}
\label{sec:ii-ter27-marcador}

La composition avec la capacité utilise l’application déclarée
$C_{s_n}\mapsto j=C_{s_n}$. Par \eqref{eq:ii-ter27-localizacion},
les niveaux sont $23+3n=3(7+n)+2$. La section exige conjointement le
résidu $2$ et un quotient au moins égal à $7$ ; le résidu isolé admettrait également
$2,5,8,\ldots$. Son projecteur et sa probabilité sont
\[
 P_{23}=\sum_{n\geq0}|23+3n\rangle\langle23+3n|,
 \qquad p_{23}=\frac{(1-r)r^{23}}{1-r^3}.
\]
Conditionner et translater $|23+3n\rangle$ vers $|n\rangle$ produit $\rho_q$.
Le poids $p_{23}$ enregistre le seuil qui disparaît de la densité normalisée.

Pour incorporer les multiplicités $d=(1,380,649)$, soit
$\iota|n,a\rangle=|23+3n\rangle\otimes|a\rangle$ dans une mémoire
de dimension $649$, avec $1\leq a\leq d_n$. C’est une isométrie.
Comprimer $\rho_r\otimes I_{649}/649$ sur son image donne la masse
\[
 p_{\mathcal D}=\frac{1-r}{649}r^{23}Z_N(q)
\]
et l’état conditionné $\rho_N(q)$. Si $\lambda_j=(1-r)r^j$,
l’intensité relative à la population de référence $\lambda_0$ est
\begin{equation}
 B_*:=\frac{\lambda_{23}+380\lambda_{26}+649\lambda_{29}}{\lambda_0}
 =10^{-23}\left(1+\frac{380}{1000}+\frac{649}{1000^2}\right)
 =1.380649\times10^{-23}.
 \label{eq:ii-ter27-intensidad}
\end{equation}
La formule conserve le spectre, le seuil, les multiplicités et la référence.
La formule historique était connue lors de la construction de ce lecteur ;
le développement opératoriel précise sa réalisation et ses conditions, au
lieu de lui attribuer rétrospectivement une prédiction indépendante.
L’interprétation comme application entre les droites d’énergie et de température
nécessite de composer les lectures dimensionnelles, objet du développement suivant.

\subsection{Courants de frontière, compression et rigidité du décalage}
\label{sec:ii-ter27-corrientes}

L’incidence d’aire déjà développée dans \ref{sec:ii-area-incidencia-explicita}
fournit $18+18$ liens tritiques. Soient $u=N_{90}$, $v=N_{120}$,
$a=u+v$ et
\[
 F=3u+4v,\qquad H_R=\epsilon_0 F,
 \quad\epsilon_0=30\Delta\tau_R>0,
 \quad\tau_R=\frac{\hbar c}{2\pi R}.
\]
L’aire physique est $\delta a\,a$, avec $\delta a$ fixée par l’incidence
et Barbero. La provenance demeure dans $u,v$ ; la même aire peut avoir
des énergies différentes. Les courants suivants utilisent cette structure
préalable sans redéfinir son hamiltonien.

Sur un lien, $T^+=|1\rangle\langle0|+|2\rangle\langle1|$ et
$T^-=(T^+)^*$ satisfont $[n,T^\pm]=\pm T^\pm$. Définissons
\[
 J_{ef}=T^+_{120,f}T^-_{90,e},\qquad
 B_c=\left(\sum_eT^+_{90,e}\right)^3
     \left(\sum_fT^-_{120,f}\right)^2.
\]
La règle $[F,AB]=[F,A]B+A[F,B]$ prouve
\[
 [F,J_{ef}]=J_{ef},\quad[a,J_{ef}]=0,
 \qquad[F,B_c]=B_c,\quad[a,B_c]=B_c.
\]
La séquence d’occupations $(0,2)\to(3,0)\to(2,1)$ a les énergies
primitives $8,9,10$ et les aires $2,3,3$. Le même saut énergétique conserve
deux actions aréales discernables.

Les coefficients $c_j=[z^j](1+z+z^2)^{18}$ comptent les occupations de
chaque canal : $c_0=1,c_1=18,c_2=171,c_3=1122$. Le caractère de frontière
\[
 P(x)=(1+x^3+x^6)^{18}(1+x^4+x^8)^{18}
\]
donne les multiplicités $g_7=324,g_8=171,g_9=1122,g_{10}=3078$.
Pour les rangs $d=(1,380,649)$, le premier triplet admissible d’énergies
consécutives commence en $8$. Les débuts inférieurs sont exclus en utilisant
$g_0=1$, $g_1=g_2=g_5=0$, $g_3=g_4=18$, $g_6=171$, $g_7=324$ et
$g_8=171$ : dans chaque triplet, au moins l’un des rangs requis fait défaut.

Une compression connexe se construit en choisissant un état $(0,2)$,
$380$ états $(3,0)$ et $649$ voisins distincts de ceux-ci par les $J_{ef}$.
L’existence d’un nombre suffisant de voisins a une preuve par comptage : chaque
occupation de somme trois possède au moins deux prédécesseurs de somme deux, et
chaque prédécesseur possède au plus $18$ successeurs. Les $380$ états
produisent donc au moins $\lceil760/18\rceil=43$ prédécesseurs
distincts et $18\cdot43=774$ voisins $(2,1)$. On peut en choisir $649$.
Le courant $B_c$ relie l’état initial à tous les $380$ états
intermédiaires ; chaque voisin final a un antécédent sélectionné. Le graphe
est connexe. Ses étiquettes et la sélection de sous-espace sont conservées.
L’isométrie basée $\iota_N$ satisfait
\begin{equation}
 H_R\iota_N=\iota_N\epsilon_0(8I+N).
 \label{eq:ii-ter27-compresion}
\end{equation}
La preuve consiste à évaluer les deux membres sur chaque vecteur de base.

Soient $P_F=\mathbf1_{N=1}$, $B_N:\mathcal D_0\to\mathcal D_1$ et
$J_N:\mathcal D_1\to\mathcal D_2$ les courants comprimés non nuls.
Pour $H_\delta=\epsilon_0(8I+N)+\delta P_F$,
\[
 [H_\delta,B_N]=(\epsilon_0+\delta)B_N,
 \qquad[H_\delta,J_N]=(\epsilon_0-\delta)J_N.
\]
Conserver les deux sauts, ou exiger leur égalité, impose $\delta=0$.
Plus généralement, une correction diagonale conserve chaque saut si ses
valeurs sont égales aux extrémités de chaque arête. La connexité la
rend constante. Cette rigidité classifie les corrections diagonales ; sa preuve
n’attribue pas cette classification à tous les opérateurs non diagonaux.

Pour $x=e^{-\epsilon_0/\tau}$ et $\tau>0$, la compression de Gibbs est
$\iota_N^*e^{-H_R/\tau}\iota_N=x^8x^N$. Sa probabilité de sélection
à partir du bord est $x^8Z_N(x)/P(x)$. Enregistrer la sélection et son complément
au moyen de $P_N\otimes|s\rangle+(I-P_N)\otimes|f\rangle$ est isométrique
et conserve l’énergie si les indicateurs sont dégénérés. Sur chaque arête
ascendante, les taux $xr_e$ et le taux inverse $r_e>0$ satisfont le bilan
détaillé pour les poids $x^N$. La connexité prouve l’unicité
stationnaire ; les masses des degrés sont $d_nx^n/Z_N(x)$. On a
ainsi spécifié une préparation thermique avec bain, non une évolution
unitaire isolée.

\Needspace{29\baselineskip}
\subsection{Équivalence thermique avec référence et échelle transportées}
\label{sec:ii-ter27-equivalencia}

\begin{theorem}[Transport de l’exposant thermique centré]
\label{thm:ii-ter27-transporte}
Soient $H_1,H_2$ des opérateurs autoadjoints finis, $\tau_i>0$ et $T$ une isométrie
d’image réductrice pour $H_2$. Les traces dans le second espace sont
prises sur cette image. Soit $\sigma_{\mathrm{ref}}$ un état normalisé
du premier espace, et soient les énergies de référence transportées
\[
 e_1=\operatorname{Tr}(\sigma_{\mathrm{ref}}H_1),\qquad
 e_2=\operatorname{Tr}(T\sigma_{\mathrm{ref}}T^*H_2).
\]
Si
\[
 \frac{H_2-e_2I}{\tau_2}T
 =T\frac{H_1-e_1I}{\tau_1},
\]
la partition centrée, l’état de Gibbs, l’entropie et
les probabilités de toutes les observables transportées sont conservés, y compris le
marqueur. Pour des états fidèles, l’égalité des états transportés
implique réciproquement cette identité avec les références ainsi définies.
\end{theorem}
\begin{proof}
Le calcul fonctionnel entrelace les exponentielles des deux exposants.
L’invariance de la trace par isométrie sur l’image donne l’égalité des
partitions centrées ; la normalisation donne $\rho_2=T\rho_1T^*$.
La cyclicité de la trace transporte les observables et le calcul spectral
transporte l’entropie. Réciproquement, prendre les logarithmes des états
fidèles produit l’égalité de $H_i/\tau_i$ à un scalaire près, égal à la
différence de leurs logarithmes de partition. Prendre l’espérance dans
$\sigma_{\mathrm{ref}}$ identifie ce scalaire à
$e_2/\tau_2-e_1/\tau_1$ ; le soustraire donne l’identité centrée énoncée.
\end{proof}

Dans le lecteur de capacité, $C=23I+3N$ et
\begin{equation}
 H_{\rm cap}=\frac{\epsilon_a\log10}{\log3}C,
 \quad\tau_{\rm clk}=\frac{\epsilon_a}{\log3},\qquad
 H_{W,\rm rel}=\epsilon_a C,
 \quad\tau_W=\frac{\epsilon_a}{\log10}.
 \label{eq:ii-ter27-hcap}
\end{equation}
Les deux exposants sont $C\log10$, bien que les énergies et les températures aient
des échelles différentes. Dans la compression du bord, la préparation
$\tau_*=\epsilon_0/\log1000$ satisfait
\[
 \frac{H_R-8\epsilon_0I}{\tau_*}\iota_N
 =\iota_N\frac{H_{\rm cap}-23\tau_{\rm clk}\log10\,I}{\tau_{\rm clk}}
 =\iota_N N\log1000.
\]
L’activité native d’horizon reste $e^{-30\Delta}$ ; cette
préparation a pour raison $\tau_*/\tau_R=30\Delta/\log1000$. Lorsqu’on fait varier
$R$, $\epsilon_0$ et $\tau_R$ changent conjointement, en conservant leur quotient.
Sans centrage, les exposants de bord et de capacité diffèrent de
$\log10\,I$, car $8\log1000-23\log10=\log10$. Leurs traces sont
$10^{-24}Z_N$ et $10^{-23}Z_N$. Le facteur $1/10$ apparaît également dans
leurs références et doit être conservé lors de la comparaison d’intensités absolues.

Si $\iota_1,\iota_2$ réalisent les mêmes rangs, degrés et marqueur,
$U=\iota_2\iota_1^*$ est unitaire entre leurs images et entrelace
le hamiltonien et le marqueur. Par le théorème précédent, tout choix de
base admissible donne le même lecteur, tous ses moments spectraux et ses
dérivées thermiques. L’unicité du représentant microscopique est
inutile pour cette conclusion. Comparer également l’aire, les courants ou
les temps de relaxation exige de transporter ces observables supplémentaires.

Enfin, pour $P_0=|\Omega_{\rm ext}\rangle\langle\Omega_{\rm ext}|$
et $\rho_N=q^N/Z_N$, on a
\[
 p_0=\operatorname{Tr}(P_0\rho_N)=Z_N^{-1},\qquad
 Z_N=1+\frac{\operatorname{Var}_{\rho_N}(P_0)}{p_0^2}.
\]
La seconde identité utilise $P_0^2=P_0$, de sorte que la variance est
$p_0(1-p_0)$. L’état normalisé et le marqueur permettent de retrouver la trace
relative ; le seuil enregistré permet de retrouver le facteur de
\eqref{eq:ii-ter27-intensidad}. Le résultat conserve ainsi la structure,
la préparation, la référence et la mémoire nécessaires à sa composition
avec la lecture d’action et de température.
